\BourbakiExercises{Exercises}

\begin{enumerate}
\def\labelenumi{\arabic{enumi})}
\tightlist
\item
  Let p be a prime number. Denote by C the set of classes of Z-modules of the form \((\mathbf{Z}/p^{2}\mathbf{Z})^{n} \oplus (\mathbf{Z}/p^{3}\mathbf{Z})^{m}\) for \(n, m \in N\) .
\end{enumerate}

\begin{enumerate}
\def\labelenumi{\alph{enumi})}
\item
  Prove that \(\mathcal{C}\) is additive but not hereditary. What is the smallest hereditary set \(\mathcal{H}\) containing \(\mathcal{C}\) ?
\item
  Prove that every short exact sequence of modules of type C splits. Describe the group \(\mathrm{K}(\mathcal{C})\) .
\item
  Construct an exact sequence
\end{enumerate}

\[
0 \to \mathrm{M} _ {0} \to \mathrm{M} _ {1} \to \mathrm{M} _ {2} \to \mathrm{M} _ {3} \to \mathrm{M} _ {4} \to 0
\]

of modules of type C such that the classes of \(M_{0} \oplus M_{2} \oplus M_{4}\) and \(M_{1} \oplus M_{3}\) in \(\mathrm{K}(\mathcal{C})\) are distinct.

\begin{enumerate}
\def\labelenumi{\alph{enumi})}
\setcounter{enumi}{3}
\tightlist
\item
  Is the natural homomorphism from \(\mathrm{K}(\mathcal{C})\) to \(\mathrm{K}(\mathcal{H})\) injective?
\end{enumerate}

\begin{enumerate}
\def\labelenumi{\arabic{enumi})}
\setcounter{enumi}{1}
\item
  Let \(\mathrm{A}\) be a ring and \(a\) and \(b\) two elements of \(\mathrm{A}\) such that \(ab = 1\) and \(ba \neq 1\) . Prove that the A-module \(\mathrm{A}(1 - ba)\) is projective and that it is stably isomorphic to the zero A-module. More precisely, prove that the A-modules \(\mathrm{A}_s\) and \(\mathrm{A}(1 - ba) \oplus \mathrm{A}_s\) are isomorphic.
\item
  Let A and B be rings. Determine the Grothendieck groups \(R(A \times B)\) and \(K_{0}(A \times B)\) of \(A \times B\) in terms of those of A and B.
\item
  Let A and B be Morita equivalent rings. Prove that the groups R(A) and R(B) are isomorphic. Do the same for \(K_{0}(A)\) and \(K_{0}(B)\) .
\item
  Let A be a ring. Denote by \(\mathcal{I}(A)\) the set of pairs \((n,p)\) with n an integer and p an idempotent in \(M_{n}(A)\) . Define a law of addition on \(\mathcal{I}(A)\) by setting \((m,p)+(n,q)=(m+n,\begin{pmatrix} p & 0 \\ 0 & q \end{pmatrix})\) .
\end{enumerate}

\begin{enumerate}
\def\labelenumi{\alph{enumi})}
\item
  Let \(\mathcal{P}(A^{\circ})\) be the monoid of classes of finitely generated projective right A-modules. Let \(\varphi: \mathcal{I}(A) \to \mathcal{P}(A^{\circ})\) be the mapping defined by \(\varphi(n,p) = \operatorname{cl}(\operatorname{Im}p)\) . Prove that \(\varphi\) is a surjective homomorphism and that elements \((m,p)\) and \((n,q)\) of \(\mathcal{I}(A)\) have the same image by \(\varphi\) if and only if there exist \(a \in M_{m,n}(A)\) and \(b \in M_{n,m}(A)\) such that ab = p and ba = q.
\item
  Deduce that the groups \(K_{0}(A)\) and \(K_{0}(A^{\circ})\) are isomorphic.
\end{enumerate}

\begin{enumerate}
\def\labelenumi{\arabic{enumi})}
\setcounter{enumi}{5}
\tightlist
\item
  Let A be a ring, G an abelian group, and T: A → G a tracial mapping (VIII, p.~115, Exercise 3). For any finitely generated projective right A-module P, denote by
\end{enumerate}

\[
\mathrm{T} _ {\mathrm{P}}: \operatorname{End} _ {\mathrm{A}} (\mathrm{P} \oplus \mathrm{A} _ {d}) \longrightarrow \mathrm{G}
\]

the unique tracial mapping that extends T (loc. cit.) and by \(\pi_{P}\) the projector of \(\operatorname{End}_{\mathrm{A}}(\mathrm{P} \oplus \mathrm{A}_{d})\) with image P and kernel \(A_{d}\) . Construct a group homomorphism \(\widetilde{\mathrm{T}} : \mathrm{K}_{0}(\mathrm{A}) \to \mathrm{G}\) such that \(\widetilde{\mathrm{T}}([\mathrm{P}]) = \mathrm{T}_{\mathrm{P}}(\pi_{\mathrm{P}})\) for every finitely generated projective right A-module P.

\begin{enumerate}
\def\labelenumi{\arabic{enumi})}
\setcounter{enumi}{6}
\tightlist
\item
  Let I be a right directed set, \((\mathrm{A}_{i}, f_{ji})\) a direct system of rings, and A its limit. Prove that \(\mathrm{K}_{0}(\mathrm{A})\) is the limit of the direct system of groups \((\mathrm{K}_{0}(\mathrm{A}_{i}), f_{ji}^{*})\) .
\end{enumerate}

¶ 8) Let A, B, and C be rings and \(f: \mathrm{A} \to \mathrm{C}\) and \(g: \mathrm{B} \to \mathrm{C}\) ring homomorphisms; suppose that \(f\) is surjective. Consider the subring D of \(\mathrm{A} \times \mathrm{B}\) consisting of the pairs \((a, b)\) such that \(f(a) = g(b)\) . Let M be an A-module, N a B-module, and \(u\) a C-isomorphism from \(\mathrm{C} \otimes_{\mathrm{A}} \mathrm{M}\) to \(\mathrm{C} \otimes_{\mathrm{B}} \mathrm{N}\) . Denote by P the D-submodule of \(\mathrm{M} \times \mathrm{N}\) consisting of the pairs \((m, n)\) such that \(u(1 \otimes m) = 1 \otimes n\) .

\begin{enumerate}
\def\labelenumi{\alph{enumi})}
\item
  Suppose that the modules M and N admit finite bases \((e_{i})_{i\in\mathrm{I}}\) and \((f_{j})_{j\in\mathrm{J}}\) and that the matrix of u in the bases \((1\otimes e_{i})\) and \((1\otimes f_{j})\) is the image of an invertible matrix \((a_{ji})\) with entries in A. Prove that the D-module P is free; more precisely, the elements \((e_{i},\sum_{j}a_{ji}f_{j})_{i\in\mathrm{I}}\) form a basis of P.
\item
  Let p and q be integers and \(X \in \mathbf{M}_{p,q}(\mathrm{C})\) and \(Y \in \mathbf{M}_{q,p}(\mathrm{C})\) matrices such that \(XY = I_{p}\) and \(YX = I_{q}\) . Prove that the square matrix \(\begin{pmatrix} X & 0 \\ 0 & Y \end{pmatrix}\) is the image of an invertible matrix with entries in A, using the identity
\end{enumerate}

\[
\left( \begin{array}{c c} X & 0 \\ 0 & Y \end{array} \right) = \left( \begin{array}{c c} I _ {p} & X \\ 0 & I _ {q} \end{array} \right) \left( \begin{array}{c c} I _ {p} & 0 \\ - Y & I _ {q} \end{array} \right) \left( \begin{array}{c c} I _ {p} & X \\ 0 & I _ {q} \end{array} \right) \left( \begin{array}{c c} 0 & - I _ {p} \\ I _ {q} & 0 \end{array} \right).
\]

\begin{enumerate}
\def\labelenumi{\alph{enumi})}
\setcounter{enumi}{2}
\item
  Suppose that the modules M and N are free and finitely generated. Deduce from a) and b) that the D-module P is projective and finitely generated.
\item
  Suppose that M and N are projective and finitely generated. Construct an A-module \(M'\) and a B-module \(N'\) such that \(M \oplus M'\) and \(N \oplus N'\) are free and finitely generated and that the C-modules \(C \otimes_{A} M'\) and \(C \otimes_{B} N'\) are isomorphic. Deduce that the D-module P is projective and finitely generated.
\item
  Prove that the canonical homomorphisms \(A \otimes_{D} P \to M\) and \(B \otimes_{D} P \to N\) deduced from the two projections are bijective (reduce as above to the case when the conditions of a) are satisfied).
\item
  Conversely, let \(P'\) be a finitely generated projective D-module; set \(M' = A \otimes_{D} P'\) and \(N' = B \otimes_{D} P'\) , and denote by \(u' : C \otimes_{A} M' \to C \otimes_{B} N'\) the canonical isomorphism. Prove that the canonical homomorphism from \(P'\) to \(M' \times N'\) induces a D-linear isomorphism from \(P'\) to the D-module associated with \((M', N', u')\) .
\end{enumerate}

\(g)\) Denote by \(\varphi_0:\mathrm{K}_0(\mathrm{D})\to \mathrm{K}_0(\mathrm{A})\times \mathrm{K}_0(\mathrm{B})\) and \(\psi_0:\mathrm{K}_0(\mathrm{A})\times \mathrm{K}_0(\mathrm{B})\to \mathrm{K}_0(\mathrm{C})\) the homomorphisms such that

\[
\varphi_ {0} ([ \mathrm{P} ]) = ([ \mathrm{A} \otimes_ {\mathrm{D}} \mathrm{P} ], [ \mathrm{B} \otimes_ {\mathrm{D}} \mathrm{P} ]) \quad \text { and } \quad \psi_ {0} ([ \mathrm{M} ], [ \mathrm{N} ]) = [ \mathrm{C} \otimes_ {\mathrm{A}} \mathrm{M} ] - [ \mathrm{C} \otimes_ {\mathrm{B}} \mathrm{N} ].
\]

Prove that the sequence

\[
\mathrm{K} _ {0} (\mathrm{D}) \xrightarrow {\varphi_ {0}} \mathrm{K} _ {0} (\mathrm{A}) \times \mathrm{K} _ {0} (\mathrm{B}) \xrightarrow {\psi_ {0}} \mathrm{K} _ {0} (\mathrm{C})
\]

is exact. Prove that if, moreover, there exists a ring homomorphism \(s: C \to A\) such that \(f \circ s = Id_{C}\) , then \(\varphi_{0}\) is injective and \(\psi_{0}\) is surjective.

\begin{enumerate}
\def\labelenumi{\arabic{enumi})}
\setcounter{enumi}{8}
\tightlist
\item
  Let A be a pseudoring, that is, an associative not necessarily unital Z-algebra, and \(\widetilde{A}\) the Z-algebra deduced from A by adjoining a unit element (III, §1, No.~2, p.~431); the set \(\widetilde{\mathrm{A}}\) is equal to \(\mathbf{Z} \times \mathrm{A}\) . Denote the homomorphism \((n, a) \mapsto n\) by \(\varepsilon : \widetilde{\mathrm{A}} \to \mathbf{Z}\) , and denote the kernel of the homomorphism \(\varepsilon^{*} : \mathrm{K}_{0}(\widetilde{\mathrm{A}}) \to \mathrm{K}_{0}(\mathbf{Z})\) associated with \(\varepsilon\) by \(\mathrm{K}_{0}(\mathrm{A})\) .
\end{enumerate}

\begin{enumerate}
\def\labelenumi{\alph{enumi})}
\item
  Prove that if A is unital, then \(\tilde{A}\) is isomorphic to the product ring \(Z \times A\) and the definition of \(K_{0}(A)\) is compatible with that given in No.~8.
\item
  Let A and B be pseudorings and \(f: \mathrm{A} \to \mathrm{B}\) a homomorphism. Let \(\widetilde{f}: \widetilde{\mathrm{A}} \to \widetilde{\mathrm{B}}\) be the ring homomorphism defined by \(\widetilde{f}((n, a)) = (n, f(a))\) . Prove that \(\widetilde{f}^*\) maps \(\mathrm{K}_0(\mathrm{A})\) into \(\mathrm{K}_0(\mathrm{B})\) . Denote the homomorphism deduced from \(\widetilde{f}^*\) by \(f^*: \mathrm{K}_0(\mathrm{A}) \to \mathrm{K}_0(\mathrm{B})\) . c) Suppose that \(f\) is surjective; denote its kernel by \(\mathfrak{a}\) and the canonical injection of \(\mathfrak{a}\) into A by \(i\) . Show that the sequence
\end{enumerate}

\[
\mathrm{K} _ {0} (\mathfrak {a}) \xrightarrow {i ^ {*}} \mathrm{K} _ {0} (\mathrm{A}) \xrightarrow {f ^ {*}} \mathrm{K} _ {0} (\mathrm{B})
\]

is exact. If, moreover, there exists a homomorphism of pseudorings \(s: B \to A\) such that \(f \circ s = Id_{B}\) , then \(i^{*}\) is injective and \(f^{*}\) is surjective (apply Exercise 8 by first taking (A, A, B) and then (D, Z, A) for the triple (A, B, C)).

\begin{enumerate}
\def\labelenumi{\alph{enumi})}
\setcounter{enumi}{3}
\tightlist
\item
  Prove that the conclusions of Exercises 7 and 8 extend to the case of pseudorings. e) Denote by \(\mathbf{M}_{\infty}(\mathrm{A})\) the pseudoring of \(\mathbf{N} \times \mathbf{N}\) matrices with entries in A having only finitely many nonzero entries. Let \(f: \mathrm{A} \to \mathbf{M}_{\infty}(\mathrm{A})\) be the mapping that sends \(a \in \mathrm{A}\) to the matrix \((b_{ij})\) with \(b_{00} = a\) and \(b_{ij} = 0\) if \((i,j) \neq (0,0)\) . Prove that \(f^{*}\) is an isomorphism from \(\mathrm{K}_0(\mathrm{A})\) to \(\mathrm{K}_0(\mathbf{M}_{\infty}(\mathrm{A}))\) (use Exercises 4 and 7).
\end{enumerate}

\begin{enumerate}
\def\labelenumi{\arabic{enumi})}
\setcounter{enumi}{9}
\tightlist
\item
  Let A be a ring. Denote by \(\mathbf{GL}_{\infty}(\mathrm{A})\) the subgroup of \(\mathbf{GL}(\mathrm{A}^{(\mathbf{N})})\) consisting of the matrices X such that X - I has only finitely many nonzero entries, and denote by \(\mathrm{E}(\mathrm{A})\) the subgroup of \(\mathbf{GL}_{\infty}(\mathrm{A})\) generated by the matrices \(I + \lambda E_{ij}\) for i, j in N with \(i \neq j\) and \(\lambda \in A\) ; it is the derived group of \(\mathbf{GL}_{\infty}(\mathrm{A})\) (II, §10, p.~422, Exercise 15). Denote the quotient group \(\mathbf{GL}_{\infty}(\mathrm{A}) / \mathrm{E}(\mathrm{A})\) by \(\mathrm{K}_{1}(\mathrm{A})\) ; its group law is written additively. Every ring homomorphism \(f : A \to B\) induces a group homomorphism \(f^{*} : K_{1}(A) \to K_{1}(B)\) .
\end{enumerate}

\begin{enumerate}
\def\labelenumi{\alph{enumi})}
\item
  For any integer \(n \geqslant 1\) , let \(\iota_{n} : \mathbf{GL}_{n}(\mathrm{A}) \to \mathrm{K}_{1}(\mathrm{A})\) be the homomorphism that sends \(X \in \mathbf{GL}_{n}(\mathrm{A})\) to the class of the matrix \(\begin{pmatrix} X & 0 \\ 0 & I \end{pmatrix}\) . Prove that if the ring A is local, then \(\iota_{n}\) is surjective for every n (use Exercise 17 of VIII, p.~42).
\item
  Suppose that the ring A is commutative. Construct a surjective homomorphism \(d_{\mathrm{A}} : \mathrm{K}_{1}(\mathrm{A}) \to \mathrm{A}^{*}\) such that \(d_{A} \circ \iota_{n} = det\) for every n.~If, moreover, A is local or Euclidean, then the homomorphisms \(d_{A}\) and \(\iota_{1}\) are inverse bijections.
\item
  Let I be a right directed set, \((\mathrm{A}_{i}, f_{ji})\) a direct system of rings, and A its limit. Prove that \(\mathrm{K}_{1}(\mathrm{A})\) is the limit of the direct system of groups \((\mathrm{K}_{1}(\mathrm{A}_{i}), f_{ji}^{*})\) .
\end{enumerate}

¶ 11) Take the assumptions and notation of Exercise 8. Denote by \(f': \mathrm{D} \to \mathrm{B}\) and \(g': \mathrm{D} \to \mathrm{A}\) the homomorphisms deduced from the two projections and by

\[
\varphi_ {1}: \mathrm{K} _ {1} (\mathrm{D}) \rightarrow \mathrm{K} _ {1} (\mathrm{A}) \times \mathrm{K} _ {1} (\mathrm{B}) \quad \text { and } \quad \psi_ {1}: \mathrm{K} _ {1} (\mathrm{A}) \times \mathrm{K} _ {1} (\mathrm{B}) \rightarrow \mathrm{K} _ {1} (\mathrm{C})
\]

the mappings defined by \(\varphi_1(z) = (g'^* (z), f'^* (z))\) and \(\psi_1(x,y) = f^* (x) - g^* (y)\) .

\begin{enumerate}
\def\labelenumi{\alph{enumi})}
\tightlist
\item
  Prove that the sequence
\end{enumerate}

\[
\mathrm{K} _ {1} (\mathrm{D}) \xrightarrow {\varphi_ {1}} \mathrm{K} _ {1} (\mathrm{A}) \times \mathrm{K} _ {1} (\mathrm{B}) \xrightarrow {\psi_ {1}} \mathrm{K} _ {1} (\mathrm{C})
\]

is exact (observe that \(f(\mathrm{E}(\mathrm{A})) = \mathrm{E}(\mathrm{C})\) ). Prove that if, moreover, there exists a ring homomorphism \(s : C \to A\) such that \(f \circ s = Id_{C}\) , then \(\psi_{1}\) is surjective.

\begin{enumerate}
\def\labelenumi{\alph{enumi})}
\setcounter{enumi}{1}
\item
  For any integer n and any element g of \(\mathbf{GL}_{n}(\mathrm{C})\) , denote by \(P_{g}\) the projective D-module associated with the modules \(A^{n}\) and \(B^{n}\) and the isomorphism from \(C \otimes_{A} A^{n}\) to \(C \otimes_{B} B^{n}\) with matrix g (Exercise 8). Construct a homomorphism \(\partial : K_{1}(\mathrm{C}) \to K_{0}(\mathrm{D})\) such that \(\partial \circ \iota_{n}(g) = [\mathrm{P}_{g}] - [\mathrm{D}_{s}^{n}]\) for every integer n and every element g of \(\mathbf{GL}_{n}(\mathrm{C})\) .
\item
  Prove that the sequence
\end{enumerate}

\[
\mathrm{K} _ {1} (\mathrm{D}) \xrightarrow {\varphi_ {1}} \mathrm{K} _ {1} (\mathrm{A}) \times \mathrm{K} _ {1} (\mathrm{B}) \xrightarrow {\psi_ {1}} \mathrm{K} _ {1} (\mathrm{C}) \xrightarrow {\partial} \mathrm{K} _ {0} (\mathrm{D}) \xrightarrow {\varphi_ {0}} \mathrm{K} _ {0} (\mathrm{A}) \times \mathrm{K} _ {0} (\mathrm{B}) \xrightarrow {\psi_ {0}} \mathrm{K} _ {0} (\mathrm{C})
\]

is exact.

\begin{enumerate}
\def\labelenumi{\arabic{enumi})}
\setcounter{enumi}{11}
\tightlist
\item
  Let \(f: \mathrm{A} \to \mathrm{B}\) be a surjective ring homomorphism, \(\mathfrak{a}\) its kernel, and \(i\) the canonical injection of \(\mathfrak{a}\) into \(\mathrm{A}\) . Construct a homomorphism \(\partial: \mathrm{K}_1(\mathrm{B}) \to \mathrm{K}_0(\mathfrak{a})\) such that the sequence
\end{enumerate}

\[
\mathrm{K} _ {1} (\mathrm{A}) \xrightarrow {f ^ {*}} \mathrm{K} _ {1} (\mathrm{B}) \xrightarrow {\partial} \mathrm{K} _ {0} (\mathfrak {a}) \xrightarrow {i ^ {*}} \mathrm{K} _ {0} (\mathrm{A}) \xrightarrow {f ^ {*}} \mathrm{K} _ {0} (\mathrm{B})
\]

is exact (argue as in Exercise 11, c)).

\begin{enumerate}
\def\labelenumi{\arabic{enumi})}
\setcounter{enumi}{12}
\item
  Denote by \(\mathcal{C}\) (resp. \(\mathcal{D}\) ) the set of classes of \(\mathbf{Z}\) -modules of finite length (resp. finitely generated \(\mathbf{Z}\) -modules). Show that the group homomorphism \(\gamma_{\mathcal{D},\mathcal{C}}\) from \(\mathrm{K}(\mathcal{C})\) to \(\mathrm{K}(\mathcal{D})\) is zero.
\item
  Let K be a commutative field of characteristic different from 2. Write
\end{enumerate}

\[
\mathrm{A} = \mathrm{K} [ \mathrm{X}, \mathrm{Y}, \mathrm{Z} ] / (\mathrm{X} ^ {2} + \mathrm{Y} ^ {2} + \mathrm{Z} ^ {2} - 1).
\]

Let M be the A-module \(\mathrm{A}^{3}/\mathrm{A}(\overline{\mathrm{X}},\overline{\mathrm{Y}},\overline{\mathrm{Z}})\) , where \(\overline{X}\) , \(\overline{Y}\) , and \(\overline{Z}\) denote the respective images of X, Y, and Z in A. Show that \(M\oplus A\) is a free A-module but that M is not a free module.

